%PIP-412035RC020
%
% Physics in Progress -- article template
%
% Options:  ec | noec     Editor's Suggestion (red border + flag)   default: noec
%           rc | oa       Rapid Communication | Original Article    default: oa
%
\documentclass[rc,ec]{physicsinprogress}

\pipvolume{2}
\pipissuedate{Feb. 2026}

% Anything else this particular article needs:
% \usepackage{siunitx}

\begin{document}

\title{It takes two: finding the normal modes in a coupled pendulum system}

\author{Brookelyn H Velmont}
\affiliation{%
Department of Physics and Astronomy, Carthage College, 2001 Alford Park Drive,
Kenosha, WI 53140, USA
}%
\pipdates{January 29, 2026}{February 5, 2026}

\begin{abstract}
The normal modes of a system of two coupled pendulums were determined by using
angular displacement data recorded with rotary motion sensors. Diagonalization
of the covariance matrix revealed two independent mode coordinates
corresponding to in-phase and out-of-phase oscillations.
\end{abstract}

\maketitle

\section{Introduction}

Synchronization is one of the most common phenomena observed in nature, yet the
physical mechanisms underlying such behavior in non-linear systems are still an
active area of study. While the motion of a single pendulum is well described by
simple harmonic motion at small angles, many real physical systems involve
multiple oscillators that exchange energy.

\section{Key Findings}

One pendulum was kept at the equilibrium position while the other was displaced
by $\theta_1$ and released from rest. To quantify correlations between the two
pendulums, the covariance matrix is constructed. Its elements are defined by
\begin{align}
A_{11} &= \left\langle (\theta_1 - \bar{\theta}_1)^2 \right\rangle, \\
A_{12} &= \left\langle (\theta_1 - \bar{\theta}_1)(\theta_2 - \bar{\theta}_2) \right\rangle,
\end{align}
where $\langle \cdot \rangle$ denotes a time average over the dataset. The
eigenvalues $\lambda$ follow from
\begin{equation}
    \det(\mathbf{A} - \lambda \mathbf{I}) = 0,
\end{equation}
and the larger eigenvalue corresponds to the dominant mode of motion. See
\cref{fig:example} for the resulting decomposition.

\begin{figure}
    \centering
    \includegraphics[width=0.35\linewidth]{Logo.pdf}
    \caption{Replace with a real figure. Referenced with \texttt{\textbackslash cref}.}
    \label{fig:example}
\end{figure}

\section{Conclusion}

The clean separation of the motion into two distinct mode coordinates
demonstrates that covariance matrix diagonalization effectively isolates the
system's normal modes.

\section*{Acknowledgments}

This research would not have been made possible without funding from the
Carthage College Physics and Astronomy Department.

\section*{Data Availability Statement}

A Jupyter notebook containing all the experimental data, data analysis, and
figure generation can be found in the supplemental
materials.\footnote{See online article posting for access to supplemental
material.}

% \bibliography{bib.bib}

\end{document}
